\documentclass[../DoAn.tex]{subfiles}
\begin{document}

\section{Đặt vấn đề}
Trong bối cảnh kỷ nguyên số và sự bùng nổ của mạng lưới Internet vạn vật (IoT), vấn đề đảm bảo an toàn thông tin trở nên cấp thiết hơn bao giờ hết. Số lượng và mức độ tinh vi của các cuộc tấn công mạng ngày càng gia tăng, đặc biệt là các hình thức tấn công từ chối dịch vụ (DoS/DDoS), quét cổng (PortScan), tấn công ứng dụng web và mạng máy tính ma (Botnet). Hậu quả của những cuộc tấn công này không chỉ gây thiệt hại lớn về kinh tế mà còn ảnh hưởng nghiêm trọng đến uy tín của các tổ chức và trải nghiệm của người dùng. Để đối phó với thực trạng đó, Hệ thống phát hiện xâm nhập (Intrusion Detection System - IDS) đóng vai trò như một tuyến phòng thủ quan trọng, giúp giám sát lưu lượng mạng và đưa ra cảnh báo kịp thời khi phát hiện các hoạt động bất thường.

Tuy nhiên, trước khối lượng dữ liệu khổng lồ và sự biến đổi không ngừng của các mã độc mới (Zero-day attacks), việc xây dựng một hệ thống IDS vừa đảm bảo độ chính xác cao, vừa có khả năng xử lý thời gian thực đang đặt ra những thách thức lớn đối với các nhà nghiên cứu và kỹ sư an toàn thông tin.

\section{Các giải pháp hiện tại và hạn chế}
Hiện nay, các giải pháp NIDS (Network Intrusion Detection System) thường được chia thành hai hướng tiếp cận chính: dựa trên luật (Signature-based IDS) và dựa trên bất thường (Anomaly-based IDS).

Hướng tiếp cận dựa trên luật, tiêu biểu là các hệ thống như Snort hay Suricata, hoạt động dựa trên cơ sở dữ liệu các dấu hiệu tấn công đã biết. Ưu điểm của phương pháp này là độ chính xác tuyệt đối đối với các cuộc tấn công đã có trong cơ sở dữ liệu và khả năng xử lý rất nhanh. Tuy nhiên, hạn chế lớn nhất là sự bất lực trước các cuộc tấn công chưa từng xuất hiện và đòi hỏi chi phí bảo trì cao để liên tục cập nhật tập luật.

Ngược lại, hướng tiếp cận dựa trên bất thường sử dụng các thuật toán Học máy (Machine Learning) và Học sâu (Deep Learning) để mô hình hóa hành vi bình thường của mạng, từ đó phát hiện các biến đổi lạ. Phương pháp này có khả năng phát hiện các cuộc tấn công mới. Dù vậy, điểm yếu cố hữu là tỷ lệ cảnh báo giả (False Positive Rate) thường ở mức cao, gây ra hiện tượng "nhiễu" cảnh báo cho các chuyên gia phân tích. Hơn nữa, các mô hình Học sâu truyền thống thường có độ trễ lớn (latency), khó áp dụng trực tiếp vào môi trường mạng lưu lượng cao.

\section\label{sec:muctieu}{Mục tiêu và định hướng giải pháp}
Nhằm khắc phục những hạn chế của các phương pháp đơn lẻ, mục tiêu chính của đồ án là nghiên cứu và xây dựng một Hệ thống phát hiện xâm nhập mạng thông minh dựa trên kiến trúc kết hợp (Hybrid IDS). Hệ thống hướng tới việc đạt được độ chính xác nhận diện cực đại (trên 99\%) và giảm thiểu tối đa tỷ lệ cảnh báo giả trong môi trường thực tế thời gian thực.

Định hướng giải pháp của đồ án bao gồm:
Thứ nhất, tích hợp hệ thống phát hiện dựa trên luật tĩnh (Snort) với bộ trích xuất đặc trưng tự động (CICFlowMeter) để nhận diện tức thì các cuộc tấn công cơ bản.
Thứ hai, ứng dụng kiến trúc mạng nơ-ron Transformer tiên tiến (cụ thể là FT-Transformer) kết hợp với kỹ thuật học quần thể (Hybrid Ensemble) để phân tích sâu các luồng dữ liệu, nhằm phát hiện các mẫu tấn công tinh vi.
Thứ ba, thiết kế và triển khai một luồng dữ liệu thời gian thực (Real-time Pipeline) cùng giao diện giám sát trực quan (Dashboard), cho phép phân tích và cảnh báo ngay lập tức các mối đe dọa.

\section{Đóng góp của đồ án}
Đồ án này có 4 đóng góp chính như sau:
\begin{enumerate}
    \item Đề xuất và triển khai kiến trúc Hybrid NIDS kết hợp giữa Signature-based (Snort) và Anomaly-based (Deep Learning), tận dụng tối đa thế mạnh của cả hai phương pháp.
    \item Ứng dụng thành công mô hình FT-Transformer kết hợp cơ chế Hybrid Ensemble để phân loại lưu lượng mạng, giải quyết bài toán mất cân bằng dữ liệu cực đoan.
    \item Xây dựng hoàn chỉnh một Data Pipeline thời gian thực: từ khâu bắt gói tin vật lý, trích xuất đặc trưng, cho đến dự đoán bằng API mô hình (FastAPI) và hiển thị kết quả lên Web Dashboard (React \& Node.js).
    \item Triển khai môi trường thực chiến (Testbed) trên hạ tầng WSL Mirrored Network, cho phép giả lập tấn công xuyên thiết bị và kiểm chứng hệ thống trong điều kiện tiệm cận thực tế.
\end{enumerate}

\section{Bố cục đồ án}
Phần còn lại của báo cáo đồ án tốt nghiệp này được tổ chức như sau:

Chương 2 trình bày ngữ cảnh của bài toán phát hiện xâm nhập mạng, điểm qua các kết quả nghiên cứu tương tự, đồng thời giới thiệu cơ sở lý thuyết về các công nghệ cốt lõi được sử dụng.

Trong Chương 3, đồ án giới thiệu kiến trúc tổng thể của hệ thống AI-Powered Hybrid NIDS đề xuất, đi sâu vào thiết kế mô hình Hybrid Ensemble cũng như luồng dữ liệu thời gian thực và hệ thống Dashboard.

Chương 4 tập trung vào phân tích đặc trưng lưu lượng mạng, các phương pháp tiền xử lý dữ liệu và thiết kế luồng xử lý phần mềm.

Chương 5 trình bày chi tiết về phương pháp thí nghiệm, các tham số đánh giá và kết quả đạt được khi huấn luyện mô hình, cũng như hiệu năng thực chiến trên hệ thống Testbed giả lập.

Cuối cùng, Chương 6 tóm tắt lại các kết quả đã đạt được, rút ra kết luận chung và đề xuất các hướng phát triển trong tương lai.

\end{document}
